\documentclass{article}
\usepackage{amsmath,amssymb}
\usepackage[vocab={calendar}]{leanify}
\newtheorem{lemma}{Lemma}
\begin{document}

\section{Notation}

We write $N(y)$ for the number of days of the year $y$.
We write $D(m,y)$ for the number of days of month $m$ of the year $y$.
We write $w(d,m,y)$ for the day of the week of the day $d$ of month $m$ of the year $y$.
We write $L(y)$ for the number of leap years up to the year $y$.

\section{Years and months}

\begin{lemma}\label{years}
$N(2000)=366$, $N(1900)=365$, $D(2,2023)=28$, $D(2,2024)=29$ and $D(4,2024)=30$.
\end{lemma}

\begin{lemma}\label{leap}
$2000$ is a leap year and $2024$ is a leap year.
\end{lemma}

\begin{lemma}\label{bound}
Let $y$ be an integer. Then $N(y)\ge365$.
\end{lemma}

\section{Weekdays}

\begin{lemma}\label{weekdays}
$w(1,1,2000)=0$, $w(25,11,2026)=4$ and $L(2000)=485$.
\end{lemma}

\end{document}
